% !PRINTED 616
\weberpage{616}
%<<CONT
eine Primzahl und der selbst, wenn er nicht die Einheitsgruppe
ist, eine transitive lineare Gruppe ist.

Dies zeigt aber unmittelbar die Darstellung der Substitutionen der Gruppe $L$ durch die Formel (5).

Wenn darin $e=1$ ist, so ist $L$ die cyklische Gruppe vom Grade $n$ mit dem Normaltheiler 1. Ist aber $e>1$, so sei $p$ eine in $e$ aufgehende Primzahl und $e=pe'$. Die lineare Gruppe $L'$ vom Grade $ne'$, die aus den Substitutionen besteht:

\[
\lambda' = (z, a_0^{ph} z + b), \quad
\begin{aligned}
h &= 0, 1, 2 \ldots e'-1\\
b &= 0, 1, 2 \ldots n-1,
\end{aligned}
\]

ist gewiss ein Theiler von $L$ vom Index $p$. Er ist aber auch normal, wie man aus der Compositionsregel (4) ohne Mühe erkennt.

Die Begriffe der transitiven linearen und der metacyklischen Gruppen decken sich also bei den Gleichungen von Primzahlgrad vollständig, und wir können daher auch in der Folge beide Ausdrücke synonym gebrauchen.

Wenn wir von einer Gruppe $P$ zu einer conjugirten Gruppe $P' = \pi^{-1}P\pi$ übergehen wollen, so geschieht dieser Uebergang dadurch, dass bei allen Permutationen von $P$ in den Cykeln die Vertauschung $\pi$ vorgenommen wird. $P'$ wird also mit $P$ bis auf die Bezeichnung der Wurzeln übereinstimmen, und ist daher, wenn $P$ linear ist, auch als linear zu bezeichnen, wenn auch eine Aenderung in der Numerirung der Wurzeln nöthig ist, um sie durch lineare Substitutionen darzustellen.

Zur Vereinfachung der Anwendung bemerke man noch, dass man die \gesperrt{volle lineare Gruppe} durch Wiederholung und Zusammensetzung der beiden Substitutionen

\weq{9}{s = (z, z+1), \quad t = (z, gz),}

und jeden Theiler $L$ der vollen linearen Gruppe ebenso aus

\weq{10}{s = (z, z+1), \quad t^{\alpha_0} = (z, a_0 z)}

ableiten kann. Man nennt daher die beiden Substitutionen (9) oder (10) die \gesperrt{erzeugenden Substitutionen dieser Gruppen}.

Setzen wir $a_0 = g^2$, so erhalten wir eine häufig vorkommende lineare Gruppe $(z, az+b)$, in der $a$ nur die quadratischen Reste von $n$ durchläuft, die von \gesperrt{Kronecker} die \gesperrt{halbmetacyklische Gruppe} genannt worden ist.

Diese Darstellung durch die erzeugenden Substitutionen gestattet einen einfachen Schluss auf die Beziehung der meta-
%CONT>>
